\section{F. Splittings and the specified Selmer problems}

\subsection{Which splittings are being computed}

\begin{statement}{F.1}{Proved}{Sections, torsors, and local conditions}
Let $k$ be a number field, $X/k$ a smooth geometrically connected
curve, and $b$ a rational point or rational tangential point.
For a characteristic quotient $U$ of the geometric fundamental
group, the section at $b$ identifies the corresponding
arithmetic extension with $U\rtimes G_k$. Continuous splittings
modulo conjugation by $U$ are the pointed set
$H^1_{\rm cts}(G_k,U)$. For a unipotent $\Q_p$-group $U$,
a Selmer scheme is a specified subset of this pointed set,
cut out by ramification and local conditions. It is not the
unrestricted set of splittings.
\end{statement}
\begin{proof}
A splitting has the form $\sigma\mapsto(c_\sigma,\sigma)$.
Multiplicativity is equivalent to
$c_{\sigma\tau}=c_\sigma\,\sigma(c_\tau)$; conjugation by
$u\in U$ changes $c_\sigma$ to $u^{-1}c_\sigma\sigma(u)$,
with the inverse convention if conjugation is defined on
the other side. These are precisely nonabelian cocycles
and their equivalence. A different reference section twists
the action and gives the corresponding bijection of
pointed sets, with no canonical choice of origin.
Unramifiedness, crystallinity at $p$, and the specified
local Kummer conditions are additional requirements on
these cocycles. They must be supplied before applying a
dimension or finiteness theorem. This is the precise
meaning used below for the repeated study of splittings
in \MS{51}{734;736}{173;178}, equations (371)--(373).
\end{proof}

\begin{statement}{F.2}{Proved}{All standard tripod dimensions through class eight}
Put $X=\mathbb P^1_\Q\setminus\{0,1,\infty\}$ and use the
unit tangent at $0$. Let $S$ be a finite set of rational
primes, $s=|S|$, and choose $p\ge5$ outside $S$.
For $U_n=\pi_1^{\rm un}(X_{\overline\Q},b)/\gamma_{n+1}$
over $\Q_p$, impose unramifiedness outside $S\cup\{p\}$
and crystallinity at $p$. Write $\operatorname{Sel}_{S,n}$
for this standard Selmer scheme, without additional refined
conditions at $S$. Then
\[
\begin{split}
 \dim\operatorname{Sel}_{S,n}
  &=2s+\sum_{\substack{3\le j\le n\\j\ {\rm odd}}}\ell_j(2),\\
 \dim H_f^1(\Q_p,U_n)&=\sum_{j=1}^n\ell_j(2),\\
 \ell_j(2)&=\frac1j\sum_{d\mid j}\mu(d)2^{j/d}.
\end{split}
\]
The first formula includes $n=1$. In particular:
\end{statement}
\begin{center}\begin{tabular}{r|rrrrr}
\toprule Class&Graded rank&Local&Global $S=0$&Global $|S|=1$&Global $|S|=2$\\\midrule
1 & 2 & 2 & 0 & 2 & 4\\
2 & 1 & 3 & 0 & 2 & 4\\
3 & 2 & 5 & 2 & 4 & 6\\
4 & 3 & 8 & 2 & 4 & 6\\
5 & 6 & 14 & 8 & 10 & 12\\
6 & 9 & 23 & 8 & 10 & 12\\
7 & 18 & 41 & 26 & 28 & 30\\
8 & 30 & 71 & 26 & 28 & 30\\
\bottomrule\end{tabular}\end{center}
For the class-$n$ free metabelian quotient $M_n$ the corresponding
exact dimensions, for $n\ge1$, are
\[
 \dim\operatorname{Sel}^{\rm met}_{S,n}
       =2s+k(k+1),\quad
 \dim H_f^1(\Q_p,M_n)=2+\frac{n(n-1)}2,\quad
 k=\left\lfloor\frac{n-1}{2}\right\rfloor.
\]
The rows at $n=1,\ldots,8$ are:
\[
\begin{array}{c|rrrrrrrr}
n&1&2&3&4&5&6&7&8\\\hline
\dim H_f^1(\Q_p,M_n)&2&3&5&8&12&17&23&30\\
\dim\operatorname{Sel}^{\rm met}_{\varnothing,n}
 &0&0&2&2&6&6&12&12
\end{array}
\]
The rows for $s=1,2$ add $2,4$ respectively. The abelian
approximation is the constant class-one row.
\begin{proof}
The central layer of $U_n$ is $\Q_p(n)^{\ell_n(2)}$.
The mixed-Tate Galois category with the stated conditions
has a graded free prounipotent Lie algebra with $s$
generators in weight one and one generator in each odd
weight at least three. Here weights are Tate degrees,
half the absolute value of the usual motivic weight.
The inputs are the Ext computations and freeness in
\citet[(2.1.3), Proposition 2.3]{DeligneGoncharov2005},
and their continuous $\ell$-adic constrained weighted realization in
\citet[Proposition 9.6, proved using Proposition 9.3 and
Theorem 9.5]{HainMatsumoto2003}.
Constrained weighted completion in their \S\S9.3--9.4 imposes
crystallinity. The degree-one group is
$\mathcal O_{\Q,S}^{\times}\otimes\Q_p$, of dimension $s$:
Kummer classes of $p$ itself are excluded by crystallinity.
The higher Ext groups are one-dimensional in odd degree
at least three and zero in even degree. The degree-two
Ext obstruction groups vanish. We use these actual
Galois completion statements, not a conjectural equality
with all of $\GRT$.

We also check the based action, because the outer
crystallinity conclusion in Hain--Matsumoto's Theorem 9.8
alone would not suffice. At the unit tangent the Belyi
normalization is $x\mapsto x^\chi$,
$y\mapsto f^{-1}y^\chi f$ with $f$ in the commutator
subgroup. It gives the unique normalized lift of the
peripheral outer action, by
\citet[Theorem 1]{HarbaterSchneps2000} and the centralizer
normalization used in B.12. At each unipotent level these
are algebraic word and conjugation formulas. Thus the
based representation factors through the same constrained
weighted completion as the outer one. This is the
additional check needed for its path-torsor cocycles.

Here is the deduction of the dimension formula, including
why it is an equality. Choose a Levi splitting and homogeneous
free generators of that Galois Lie algebra. A cocycle to $U_n$
can be uniquely normalized on the Levi $\Gm$, since all
weights of $\Lie U_n$ are strictly negative. It then
amounts to a graded Lie section into
$\Lie U_n\rtimes\mathfrak u_{\rm Gal}$. On a free generator
of degree $j$, its value may be the given generator plus
an arbitrary element of $(\Lie U_n)_j$; freeness supplies
the unique homomorphism for all such choices. Conversely
every normalized cocycle has precisely these values.
Thus the scheme is, noncanonically, an affine space with
two coordinates for each degree-one Galois generator and
$\ell_j(2)$ for its generator in odd degree $j\ge3$.
The choices of generators and Levi splitting explain the
noncanonicity; the dimension and the Selmer functor do not
depend on them. This also proves compatibility of the
formulas with the tower and describes its inverse-limit
Selmer functor as the corresponding product of coordinates.

At $p$, the crystalline comparison identifies
$H_f^1(\Q_p,U_n)$ with $U_n^{\rm dR}/F^0$.
All Hodge degrees of this tripod Lie algebra are negative,
so $F^0$ is trivial. Witt's formula proves the local
dimension. For the free metabelian Lie algebra on two
generators the derived ideal is the free rank-one module
generated by $[x,y]$ under the commuting operations
$\ad x,\ad y$. Consequently its degree-$j$ rank is $j-1$
for $j\ge2$. Substitution into the same cocycle argument
gives both metabelian formulas and their full inverse
tower. The code computes every row by exact integer
arithmetic; it does not enumerate rational or integral
points.
\end{proof}

\begin{statement}{F.3}{Known}{The first dimension drops and integral finiteness}
For each of $U_n$ and $M_n$, the first strict inequality
\[
 \dim\operatorname{Sel}_{S,n}<\dim H_f^1(\Q_p,U_n)
\]
(with $M_n$ substituted in its case) occurs at class
$1,2,4$ for $s=0,1,2$, respectively. Each inequality
implies that $X(\Z[1/S])$ is finite. It does not identify
that set with the corresponding Selmer locus.
\end{statement}
\begin{proof}
Direct comparison of the exact rows of F.2 proves the
minimality of all three thresholds. The abelian
approximation alone gives the strict inequality only
for $s=0$. For the nonabelian towers,
\citet[Theorem 1 and \S3]{Kim2009} supplies Zariski density
of the unipotent Albanese map on each residue disk and
the dimension criterion for finiteness. His notation
$U_2$ begins with abelianization; our $U_1$ is his $U_2$.
We make this indexing shift explicit before invoking
his theorem.

The global-to-local localization morphism has image
closure of dimension at most the global Selmer dimension.
Under the strict inequality there is a nonzero regular
function on the local space vanishing on that closure.
Its pullback by the unipotent Albanese map is a nonzero
Coleman function on each residue disk. The map for the
metabelian quotient is also Zariski dense: it is the
composition of the dense full map and the surjective
quotient morphism. There are only finitely many integral
residue disks at $p$. On each compact disk a nonzero
analytic function has finitely many zeros, so the
global integral points, whose Kummer classes lie in
the localization image, form a finite set. This proves
the manuscript's proposed mechanism of refining
splittings in this explicitly defined tower; it is not
a finiteness theorem for unrestricted sections.
\end{proof}

\begin{statement}{F.4}{Refuted}{Abelian approximations do not separate every cusp}
For the tripod the unrestricted abelian section sets are
\[
 H^1(G_\Q,\widehat\Z(1)^2)
       =\left(\varprojlim_m\Q^\times/(\Q^\times)^m\right)^2,
 \qquad
 H^1(G_\Q,\Q_p(1)^2)
       =\left(\widehat{\Q^\times}^{\,p}\otimes_{\Z_p}\Q_p\right)^2.
\]
They are infinite. The two tangential sections at $0$ and $1$,
using respectively parameters $z$ and $1-z$ with unit tangent,
have the same class in both sets. Thus the assertion that
distinct cusps remain nonconjugate in every abelian
approximation is false. This counterexample concerns an
approximation, and does not identify the full profinite
cuspidal sections.
\end{statement}
\begin{proof}
The finite Kummer exact sequence and Hilbert 90 give
$H^1(G_\Q,\mu_m)=\Q^\times/(\Q^\times)^m$.
Here $\widehat{\Q^\times}^{\,p}$ denotes the power
completion $\varprojlim_j\Q^\times/(\Q^\times)^{p^j}$.
For this infinite-rank group it is not the completion
over all finite $p$-group quotients.
The inverse systems of finite invariant root groups
are Mittag--Leffler, so the $\varprojlim^1 H^0$ term
vanishes and passage to $\widehat\Z(1)$ gives the first
formula. The identical argument for powers of $p$ gives
the integral $p$-adic formula. A continuous cocycle from
compact $G_\Q$ to $\Q_p(1)$ has bounded image, hence after
one common multiplication by a power of $p$ it is
$\Z_p(1)$-valued. This proves the tensor assertion.

These groups are quite explicit. Write
$\Q^\times=\{\pm1\}\oplus\bigoplus_q\Z$ using valuations
at rational primes $q$. In the profinite formula an
element is a sign and a family $a_q\in\widehat\Z$ such
that, for every positive $m$, $a_q$ is zero modulo $m$
at all but finitely many $q$. In the $p$-adic integral
formula the analogous family lies in $\Z_p$ and tends
to zero $p$-adically outside finite sets, together with
the sign torsion if $p=2$. Tensoring permits a common
bounded denominator. This computes the named abelian
cohomology groups, including their possible infinite
ramification; it does not mistake a direct sum for
its completion.

The two coordinates of abelian path torsors are the
Kummer torsors of $z$ and $1-z$. At a tangential
endpoint, use the leading coefficient of a vanishing
function in the chosen parameter. At the unit tangent
of $0$ these two regularized values are $(1,1)$;
at the unit tangent of $1$ they are again $(1,1)$.
All their Kummer classes vanish. Hence both section
classes are the origin relative to the first section.
Equivalently the connecting path has zero abelian
Galois cocycle. The corresponding standard Selmer
classes for $S=\varnothing$ also coincide, consistently
with the dimension zero in F.2.
\end{proof}

\begin{statement}{F.5}{Proved}{The effect of the degree-seven nonpentagonal element}
Let $\ell\ge5$ and let $\beta$ be the pro-$\ell$ element
of B.12 with normalized conjugator $f$ having leading
Lie term $q_7$. On $U_n$ it is the identity for $n\le7$.
In class eight, with the normalization
$\beta(x)=x,\ \beta(y)=f^{-1}yf$, its first change is
\[
 \beta(\log x)=\log x,\qquad
 \beta(\log y)=\log y+[\log y,q_7]\pmod{\text{degree}>8}.
\]
It transports cocycles for $\rho$ bijectively to cocycles
for $\rho^\beta=\beta\rho\beta^{-1}$. It need not act
on the section set for the \emph{fixed} action $\rho$.
All tripod finiteness thresholds in F.3 occur before
this change can be detected.
\end{statement}
\begin{proof}
Conjugation by $\exp(q_7+\text{higher terms})$ changes
a degree-one generator first in degree eight, and its
first term is the displayed bracket. This proves the
identity through class seven and the explicit class-eight
action; it also descends to every characteristic quotient.
If $c_{\sigma\tau}=c_\sigma\rho(\sigma)c_\tau$, application
of $\beta$ gives the cocycle equation for $\rho^\beta$.
Coboundaries are transported in the same way, proving
the bijection. For a fixed action, the same calculation
requires $\beta\rho(\sigma)=\rho(\sigma)\beta$, or an
additional specified intertwining cocycle; the two
cycle equations alone do not give that condition.
Selmer conditions are preserved only when their local
structures are transported as well. In particular a
non-arithmetic element cannot be used to change a count
of fixed-action Selmer classes by this argument.
At classes $1,2,4$ it acts identically, so no such
change is present at the computed dimension-drop levels.
\end{proof}

\subsection{Part F(c). Two punctured elliptic curves, with exact local dimensions}

This subsection implements the successive examination of splittings proposed
in the closing pages of \S51, folio 734, scan 173, and folio 736, scan 178.
It concerns the geometric group of type $(1,1)$ occurring in \S52,
folios 743--746, scans 191--197. It supplies calculations for the abelian,
nilpotent and free metabelian unipotent approximations, with their Galois
actions. Splittings always mean conjugacy classes of splittings satisfying
the specified local conditions. Their parameter schemes can have positive
dimension even when the set of integral points is finite.

Let $E_d:y^2=x^3-d^2x$, $X_d=E_d\setminus\{O\}$, and take $d=1,6$.
We use $p=5$ and the rational base point $b=(0,0)$. Let $\mathcal E_d$
be the minimal proper regular model over $\mathbb Z$ and let
$\mathcal X_d=\mathcal E_d\setminus\overline{\{O\}}$. This model contains
the integral affine points on the displayed Weierstrass equation. We allow
a finite set $S$ of denominator primes, with $5\notin S$, and put $s=|S|$.
The bad primes are $B_1=\{2\}$ and $B_6=\{2,3\}$; put
$T=S\cup B_d\cup\{5\}$. They are not automatically counted in $s$:
at a bad prime outside $S$ the integral local condition is retained.

Write $U_n=\pi_1^{\rm un}(X_{d,\overline{\mathbb Q}},b)/\gamma_{n+1}$,
over $\mathbb Q_5$, and $V=T_5E_d\otimes\mathbb Q_5$. Thus $U_1=V$.
Write $M_n$ for the corresponding class-$n$ quotient of the free metabelian
quotient $U/[U',U']$. Both towers have $L_1=V$, $L_2=\mathbb Q_5(1)$.
For $U$ the higher central layers are $L_n=\operatorname{Lie}_n(V)$;
for $M$ write $L_n^{\rm met}$.
For each finite reduction type $\Sigma$, impose at $\ell\ne5$ the
Zariski closure of the local Kummer image of the indicated integral
component, or of the indicated punctured component at $\ell\in S$.
At $5$ impose crystallinity. Denote the resulting finite union of
refined Selmer schemes by $\operatorname{Sel}_{S,n}$, or
$\operatorname{Sel}^{\rm met}_{S,n}$. Its dimension means the maximum
dimension of its nonempty components. This convention permits the cuspidal
local component at every prime in $S$, even when it has no global integral
point. These schemes still contain the Kummer images of all
$\mathcal X_d(\mathbb Z[1/S])$.

\begin{statement}{F.6}{Proved}{The examples have ranks zero and one, and the usual Selmer ranks agree}
The curves $E_1$ and $E_6$ have Mordell--Weil ranks $0$ and $1$, respectively.
They both have good ordinary reduction at $5$, with $a_5=-2$.
Their Tate--Shafarevich groups are finite; consequently
\[
 \dim H_f^1(\mathbb Q,V)=r,\qquad r=0\ \hbox{or}\ 1,
 \qquad \dim H_f^1(\mathbb Q_5,V)=1.
\]
Here $H_f^1$ uses the Bloch--Kato local conditions at every place.
\end{statement}
\begin{proof}
The pure $2$-Selmer rank of a squarefree congruent-number curve is the
nullity of its Monsky matrix. In the additive Legendre-symbol convention,
for $d=2^\epsilon P_1\cdots P_t$ the even matrix is
\[
 M_e=\begin{pmatrix}D_2&A+D_2\\ A^{\mathsf T}+D_2&D_{-1}\end{pmatrix},
 \qquad
 A_{ij}=\left[\frac{P_j}{P_i}\right]\ (i\ne j),\quad
 A_{ii}=\sum_{j\ne i}A_{ij},
 \quad (D_h)_{ii}=\left[\frac h{P_i}\right].
\]
The empty odd matrix for $d=1$ has nullity zero. For $d=6$, $t=1$,
$P_1=3$, $A=0$, $D_2=D_{-1}=(1)$, so
$M_e=\left(\begin{smallmatrix}1&1\\1&1\end{smallmatrix}\right)$ has
nullity one. Thus the ranks are at most zero and one. This is the
explicit descent formula in \citet[\S2, (2.1)--(2.3)]{FE_DasSaikia2024},
attributed there to Monsky's appendix to Heath--Brown.

For $E_6$, $P=(12,36)$ satisfies
\[
 2P=(25/4,-35/8),\qquad
 4P=(1442401/19600,1726556399/2744000),\qquad 8P\ne O.
\]
Direct counts give $|E_6(\mathbb F_5)|=|E_6(\mathbb F_7)|=8$.
The kernel of reduction at a good prime is a pro-prime formal group, so
reduction is injective on torsion of order prime to that prime. Combining
the two primes shows that the order of rational torsion divides $8$.
Thus $P$ is not torsion, and the rank is one. The same counts at $3,5$
for $E_1$ are $4,8$; its four visible rational $2$-torsion points therefore
exhaust $E_1(\mathbb Q)$. The counts at $5$ also give $a_5=6-8=-2$ for
both curves, proving ordinarity.

The $2$-descent alone does not prove finiteness of the $5$-primary
Tate--Shafarevich group. We use a separate verified input. The primary
database entries $32.a3$ and $576.c3$ record analytic ranks zero and one
\citep{FE_LMFDBE1,FE_LMFDBE6}. Their conductors are $32$ and $576$.
The database's reliability statement specifies that analytic ranks at
most three at conductor at most $500000$ are rigorous modular-symbol
computations \citep{FE_LMFDBReliability}. Kolyvagin's own
\citet[Theorem 1, printed p.~4; also Example 1, p.~11]{FE_Kolyvagin1990}
states that, for a modular elliptic curve of analytic rank at most one,
the algebraic rank agrees and the Tate--Shafarevich group is finite.
The hypotheses hold here. The Kummer exact sequence after tensoring
with $\mathbb Q_5$ gives the global assertion; the formal logarithm gives
the one-dimensional local assertion.

The program \texttt{code/f\_elliptic\_dimensions.py} checks both descent
matrices, all displayed point operations, and all four reduction counts
by exact rational or finite-field arithmetic. The analytic-rank input is
identified separately in its output and is not purported to be recomputed.
\end{proof}

\begin{statement}{F.7}{Proved}{Full and metabelian local dimensions through class eight}
For either curve the exact local dimensions are
\[
 \begin{array}{c|rrrrrrrr}
 n&1&2&3&4&5&6&7&8\\\hline
 \dim H_f^1(\mathbb Q_5,U_n)&1&2&4&7&13&22&40&70\\
 \dim H_f^1(\mathbb Q_5,M_n)&1&2&4&7&11&16&22&29
 \end{array}
\]
The complete linear data used below are
\[
 \begin{array}{c|rrrrrrrr}
 n&1&2&3&4&5&6&7&8\\\hline
 w_n=\dim L_n&2&1&2&3&6&9&18&30\\
 \operatorname{tr}(c\mid L_n)&0&-1&0&-1&0&-1&0&-2\\
 \dim L_n^{c=-1}&1&1&1&2&3&5&9&16\\
 \dim L_n^{\rm met}&2&1&2&3&4&5&6&7\\
 \dim (L_n^{\rm met})^{c=-1}&1&1&1&2&2&3&3&4
 \end{array}.
\]
The degree is the lower-central degree, equal here to the negative
geometric weight. It is not the doubled Tate grading of the tripod.
\end{statement}
\begin{proof}
The topological fundamental group of a once-punctured torus is free on
two generators. The central graded Lie algebra of its unipotent
completion is therefore the free Lie algebra on $V$. The Hodge
filtration on $V$ has one generator in degree $0$ and one in degree
$-1$. A Lie word of length at least two in just the degree-zero
generator vanishes. Hence $F^0L_n=0$ for $n\ge2$, whereas
$\dim F^0V=1$. The de Rham comparison therefore gives
\[
 \dim(U_n^{\rm dR}/F^0)=\sum_{m=1}^n w_m-1,
 \qquad w_m=\frac1m\sum_{e\mid m}\mu(e)2^{m/e}.
\]
There are no Frobenius-fixed vectors in the pure negative-weight
layers. The crystalline torsor classification identifies the displayed
de Rham quotient with $H_f^1(\mathbb Q_5,U_n)$; alternatively its
dimension follows by successive central extensions and
\citet[Corollary 3.1.11]{FE_Betts2023}.

Complex conjugation on $V$ has eigenvalues $1,-1$. The character form
of the Witt formula is
\[
 \operatorname{tr}(c\mid L_m)
 =\frac1m\sum_{e\mid m}\mu(e)
       \bigl(\operatorname{tr}(c^e\mid V)\bigr)^{m/e},
 \qquad
 \dim L_m^{c=-1}=\frac{w_m-\operatorname{tr}(c\mid L_m)}2.
\]
For the free metabelian Lie algebra the elements
$\operatorname{ad}_x^a\operatorname{ad}_y^b[x,y]$, $a+b=m-2$,
form a basis in degree $m\ge2$. Thus the dimension is $m-1$ and the
negative conjugation eigenspace has dimension $\lfloor m/2\rfloor$.
The Hodge argument is unchanged.

As an independent finite check, the program exhausts all binary words
of each length $1$ through $8$, selects the Lyndon words, and counts
them and their signs $(-1)^{\#y}$. It compares these counts with the
Witt formulas and with the explicit metabelian basis. It also verifies
every CM bidegree multiplicity. This is exhaustive at the stated word
lengths, and does not sample Galois cohomology.
\end{proof}

\begin{statement}{F.8}{Proved}{Exact low-level global dimensions and unconditional finiteness}
For either curve, and every finite $S$ with $5\notin S$,
\[
 \dim\operatorname{Sel}_{S,1}=r,
 \qquad \dim\operatorname{Sel}_{S,2}=r+s.
\]
The same identities hold for the metabelian tower. For $E_1$ the first
strict global--local dimension inequality is at class one, for all
$s=0,1,2$. For $E_6$ and $s=0$ it is at class two. These inequalities
prove finiteness of the corresponding integral sets by the unipotent
Albanese method. In particular
$X_1(\mathbb Z[1/S])=\{(-1,0),(0,0),(1,0)\}$ for every $S$.
\end{statement}
\begin{proof}
At every $\ell\ne5$, potential good reduction gives
$H^1(\mathbb Q_\ell,V)=0$: $H^0$ and $H^2$ vanish by weights and local
duality, and the local Euler characteristic is zero. The ordinary
abelian Selmer condition is thus independent of $S$, and F.6 proves
the first identity.

For class two, the involution $[-1]$ fixes $b$ and commutes with Galois.
It acts by $-1$ on $V$ and $+1$ on the central line. Its eigenspaces
therefore give a Galois-equivariant vector-space splitting
\[
 \operatorname{Lie}(U_2)=V\oplus\mathbb Q_5(1),
 \qquad [(v,z),(v',z')]=(0,v\wedge v').
\]
In these coordinates a cocycle consists of $(a,b)$ with
$da=0$ and $db=-\tfrac12a\mathbin{\cup}a$.
For rank zero only $a=0$ occurs. For rank one take the cocycle of the
rational point $P=(12,36)$ and scale its coordinates by
$(a,b)\mapsto(ta,t^2b)$. This supplies an algebraic lift for every
$t\in\mathbb Q_5$ in the full abelian Selmer line; it remains
crystalline at $5$, since $P$ is integral there and the scaling is an
endomorphism of the class-two Galois group object. Consequently the
global class-two fibres before imposing the remaining local conditions
are torsors under
\[
 H^1(G_{\mathbb Q,T},\mathbb Q_5(1))
   =\mathbb Z[1/T]^\times\otimes\mathbb Q_5.
\]
There is no stabilizer quotient: the relevant twisted $H^0$ groups
vanish, since their central graded pieces have negative weights.

The valuation map from this unit space to
$\bigoplus_{\ell\in T}\mathbb Q_5$ is an isomorphism. Crystallinity
at $5$ fixes the $5$-valuation. At $\ell\in B_d\setminus S$, the
integral local Kummer images are finite, and fixing a reduction type
fixes the corresponding central valuation. At a prime in $S$ the
cuspidal component has Zariski closure the full affine central line:
the puncture is rational, and its leading commutator is the nonzero
generator of $\mathbb Q_5(1)$, so its unbounded integral contact orders
are Zariski dense in that line. Thus the type which is cuspidal at
every prime in $S$ leaves exactly $s$ unit parameters. Types which
are noncuspidal at some of them leave fewer parameters. The unit
valuation isomorphism both satisfies the required constants at bad
primes and proves that the maximal dimension is exactly $r+s$.
The finiteness of the non-$5$ integral local images is
\citet[Theorem 0.1 and Corollary 0.2]{FE_KimTamagawa2008}.

F.7 gives the local dimensions $1,2$ at these levels. A proper
algebraic closed subset of the local period space is cut out by a
nonzero regular function. The pullback is a Coleman function; the
unipotent Albanese map is Zariski dense on each residue disk, so it
does not vanish identically there. On the compact integral disks a
nonzero such analytic function has finitely many zeroes. There are
finitely many disks and finitely many reduction types. This proves
finiteness. The original density theorem is
\citet[Theorem 1, p.~93; finiteness deduction in \S3]{FE_Kim2009}.
His indexing starts with $U_2$ abelian, so his $U_{n+1}$ is our $U_n$.
The elliptic rank-at-most-one application appears explicitly in
\citet[Corollary 0.3 and its proof in \S2]{FE_KimTamagawa2008}, whose
$U_3$ is our class two. The modern refined formulation also follows
from \citet[Theorem A, Remark 1.4, and Lemma 3.2]{FE_LLM2023}.

For rank one the abelian localization is nonzero: a non-torsion
rational point has nonzero elliptic $5$-adic logarithm. Thus class one
does not give a dimension inequality. At $s=0$ class two does. The
description of $E_1(\mathbb Q)$ in F.6 gives the last assertion.
\end{proof}

\begin{statement}{F.9}{Open}{Higher elliptic dimensions: exact defects and a conditional complete table}
For $n\ge3$ define the actual finite-dimensional arithmetic invariants
\[
 \delta_n(E)=\dim H_f^1\bigl(\mathbb Q,L_n^*(1)\bigr),\qquad
 \delta_n^{\rm met}(E)=
       \dim H_f^1\bigl(\mathbb Q,(L_n^{\rm met})^*(1)\bigr).
\]
The exact strict global graded dimensions are
\[
 \dim H_f^1(\mathbb Q,L_n)=\dim L_n^{c=-1}+\delta_n(E),
 \quad
 \dim H_f^1(\mathbb Q,L_n^{\rm met})
       =\lfloor n/2\rfloor+\delta_n^{\rm met}(E).
\]
Consequently for $n\ge2$ there are the unconditional bounds
\[
 \begin{split}
 \dim\operatorname{Sel}_{S,n}
   &\le r+s+\sum_{m=3}^n\bigl(\dim L_m^{c=-1}+\delta_m(E)\bigr),\\
 \dim\operatorname{Sel}^{\rm met}_{S,n}
   &\le r+s+\sum_{m=3}^n
           \bigl(\lfloor m/2\rfloor+\delta_m^{\rm met}(E)\bigr).
 \end{split}
\]
If the relevant defects vanish through the displayed level, these
bounds are equalities. In particular, subject to that explicitly
stated hypothesis, the exact global dimensions through class eight are
\[
\begin{array}{cc|rrrrrrrr}
 \multicolumn{2}{c|}{\text{full free tower}}&1&2&3&4&5&6&7&8\\
 r&s&&&&&&&&\\\hline
 0&0&0&0&1&3&6&11&20&36\\
 0&1&0&1&2&4&7&12&21&37\\
 0&2&0&2&3&5&8&13&22&38\\
 1&0&1&1&2&4&7&12&21&37\\
 1&1&1&2&3&5&8&13&22&38\\
 1&2&1&3&4&6&9&14&23&39
\end{array}
\]
\[
\begin{array}{cc|rrrrrrrr}
 \multicolumn{2}{c|}{\text{metabelian tower}}&1&2&3&4&5&6&7&8\\
 r&s&&&&&&&&\\\hline
 0&0&0&0&1&3&5&8&11&15\\
 0&1&0&1&2&4&6&9&12&16\\
 0&2&0&2&3&5&7&10&13&17\\
 1&0&1&1&2&4&6&9&12&16\\
 1&1&1&2&3&5&7&10&13&17\\
 1&2&1&3&4&6&8&11&14&18
\end{array}
\]
The first strict dimension inequalities for rank one are at classes
$2,3,4$ for $s=0,1,2$. The class-two conclusion is unconditional;
the class-three conclusion uses $\delta_3(E_6)=0$, and the class-four
conclusion uses $\delta_3(E_6)=\delta_4(E_6)=0$. The full and metabelian
groups agree through class four. This audit does not establish those
vanishings, or the defects through class eight; the corresponding
unconditional numerical higher Selmer dimensions remain unresolved here.
\end{statement}
\begin{proof}[Proof of the reductions and conditional assertions]
For a geometric representation $W$, Poitou--Tate duality gives
\[
 h_f^1(\mathbb Q,W)=h^0(\mathbb Q,W)
 +h_f^1(\mathbb Q,W^*(1))-h^0(\mathbb Q,W^*(1))
 +h_f^1(\mathbb Q_5,W)-\dim W^{c=1}.
\]
This precise formula is recorded in
\citet[(5.6)]{FE_LLM2023}. For $W=L_n$, $n\ge3$, both global
$H^0$ terms vanish by weights, and $h_f^1(\mathbb Q_5,W)=\dim W$
by F.7. This proves the exact graded formula, including its defect.
For $L_2=\mathbb Q_5(1)$ the dual invariant term is one and the strict
global group is zero; it cannot be treated by the $n\ge3$ formula.
Allowing denominator primes supplies the class-two unit term $s$ of F.8.

Here potential good reduction is useful beyond dimension estimates.
For $n\ge3$ and $\ell\ne5$,
\[
 H^i(\mathbb Q_\ell,L_n)=0\quad(i=0,1,2).
\]
Indeed after a finite extension the representations are unramified and
pure of weight $-n$; the dual twist has weight $n-2$. Neither weight
is zero. Thus the two invariant groups vanish, and local duality and
Euler characteristic give the assertion. The same argument works for
the metabelian layers. At $5$, $F^0L_n=0$ and the dual twist has
$F^0$ equal to its whole de Rham space, with no invariant vectors.
It follows that
\[
 H^1(\mathbb Q_5,L_n)=H_f^1(\mathbb Q_5,L_n),\qquad
 H_f^1(\mathbb Q_5,L_n^*(1))=0.
\]
Poitou--Tate therefore identifies
\[
 H^2(G_{\mathbb Q,T},L_n)^*
       =H_f^1(\mathbb Q,L_n^*(1)).
\]
There is no dependence on the extra finite primes of $T$ in these
groups: the displayed local groups at such primes vanish.

In a central extension
$1\to L_n\to U_n\to U_{n-1}\to1$, the obstruction to lifting a
Galois cocycle lies in this $H^2$, and every nonempty fibre is a
torsor under $H^1(G_{\mathbb Q,T},L_n)$, since the twisted $H^0$
stabilizers vanish. This proves the unconditional upper bound by
induction. If $\delta_n=0$, all lower global classes lift. At a
non-$5$ place their lift is unique. At $5$ the crystalline condition
also lifts: the de Rham quotients
$U_n^{\rm dR}/F^0\to U_{n-1}^{\rm dR}/F^0$ are surjective by
strictness of the Hodge filtration, and the whole central $H^1$ is
crystalline. Thus it is the fibre of the local lift.
Consequently no additional local equation appears in degree $n$.
The map of refined Selmer schemes is then a torsor, locally on the
base, under the vector group of dimension $\dim L_n^{c=-1}$.
Starting with the maximal class-two component of dimension $r+s$
from F.8 proves equality, rather than merely a conditional upper
bound. The identical argument proves the metabelian assertion.
Summing the rows of F.7 gives the two tables.

For completeness the remaining groups have an explicit CM description.
Over $K=\mathbb Q(i)$ write $V=\psi\oplus\overline\psi$ and
$\psi\overline\psi=\chi_5$. The multiplicity of
$\psi^a\overline\psi^b$ in $L_{a+b}$ is
\[
 \ell_{a,b}=\frac1{a+b}
   \sum_{e\mid\gcd(a,b)}\mu(e)
                  \binom{(a+b)/e}{a/e},\qquad a,b\ge1.
\]
It is one in the metabelian layer. The dual-twist characters are
$\psi^{1-a}\overline\psi^{1-b}$, with these multiplicities;
restriction to $K$ and the order-two descent involution recover the
groups over $\mathbb Q$. The program lists every pair $(a,b)$ through
$a+b=8$. This reduces the missing arithmetic to a finite, named list
of Hecke-character Selmer groups, but Hodge weights alone do not
force their vanishing. In particular $L_3=V(1)$, so already
$\delta_3(E)=h_f^1(\mathbb Q,V^*)$ is a distinct problem from
the Mordell--Weil rank calculation.

There are meaningful alternative theorems, but they do not certify the
uncomputed rows. \citet[Definition 2.2 and Corollary 2.9]{FE_Bars2007}
proves the requisite second-cohomology vanishing under geometric
regularity of $5$ for $E$ and $K(E[5])$; that regularity has not been
verified here. \citet[Theorem 0.1]{FE_Kim2010CM} proves eventual
global--local dimension inequality unconditionally for a smaller CM
quotient $W$; his Theorem 0.2 gives an explicit bound subject to
specified nonvanishing values of $p$-adic $L$-functions. The quotient
$W$ kills the Lie bidegrees with both entries at least two and is a
quotient of the metabelian group. Thus it gives an alternative
unconditional eventual finiteness argument for both curves, after
enlarging $S$ to include bad primes if necessary. Its existence theorem
does not supply a first level at most eight or the full-table defects.

Any conditional strict inequality established above implies finiteness
by the same density and analytic argument as F.8; the general effective
refinement is \citet[Theorem 6.2.1(A),(B), and \S6.3]{FE_Betts2023}.
For transfer to LM6, the lemma must retain the specified local
conditions, the rank or finite-$\Sha$ input, the actual Galois action,
and every needed $\delta_m=0$ hypothesis. Replacing these hypotheses
by an assertion about a normalizer Lie algebra supplies no such
arithmetic vanishing theorem.
\end{proof}


\subsection{Exact lifting statements and the remaining section problems}

\begin{statement}{F.10}{Proved}{The obstruction at each successive central extension}
Let $1\to A\to U\to V\to1$ be a central extension of
topological $G_k$-groups, with $A$ abelian and with
continuous local lifts for the cocycles under consideration.
For a $V$-cocycle $c$, choose a continuous lift $\widetilde c$.
Its defect
\[
 o_c(\sigma,\tau)=
 \widetilde c_\sigma\,\sigma(\widetilde c_\tau)
                    \widetilde c_{\sigma\tau}^{-1}\in A
\]
defines a class in $H^2_{\rm cts}(G_k,A)$ independent
of the lift. The class vanishes exactly when $c$ lifts.
For a fixed representative $c$, its lifts modulo
$A$-conjugacy form a torsor under $H^1_{\rm cts}(G_k,A)$.
Passing to $U$-conjugacy further divides by the stabilizer
of $c$ in $V$. Local conditions impose the corresponding
local obstruction and localization conditions.
\end{statement}
\begin{proof}
Associativity and the cocycle equation in $V$ give the
usual degree-two cocycle identity for $o_c$. Centrality
ensures that no additional conjugation on $A$ occurs.
Replacing $\widetilde c$ by $a_\sigma\widetilde c_\sigma$
changes its defect by the coboundary of the $A$-cochain
$a$. Vanishing therefore permits correction to an
actual cocycle, and the quotient of two corrected
lifts is an $A$-cocycle. Quotienting by conjugation gives
the stated torsor and stabilizer action.
For unipotent groups over $\Q_p$, exponential coordinates
supply the required continuous lifts. For profinite
extensions this local-lift hypothesis must be checked,
or the assertion formulated using continuous extensions
and their nonabelian cohomology instead.
Applying the construction at every lower-central stage
gives the exact successive test proposed at scans173
and178. Merely naming its $H^1$ and $H^2$ groups does
not evaluate them.
\end{proof}

\begin{statement}{F.11}{Proved}{The condition for restricting a section to a cover}
Suppose $1\to N\to P\to H\to1$ is an exact geometric
finite-cover sequence, stable under the arithmetic
action after a specified finite extension of the base
field. If the chosen approximation preserves this
exact sequence, a section of its arithmetic extension
restricts to the cover exactly when its induced section
of $H\rtimes G_k$ is the section specifying that cover,
up to the stated $H$-conjugacy. Restriction after an
open subgroup of $G_k$ is possible whenever the
corresponding finite cocycle becomes trivial there.
\end{statement}
\begin{proof}
The arithmetic group of a fixed connected cover is
the inverse image of its finite section in the
quotient extension. A section lands in that inverse
image precisely under the displayed condition.
Changing the finite section by conjugation changes
the chosen lift of the base point and yields exactly
the stated ambiguity. A continuous finite cocycle
and the finite action factor through a finite
quotient of $G_k$; their common open kernel kills
that cocycle. Exactness then identifies the geometric
kernel with the approximation of $N$, so the
restricted section is a section for the cover.
If that comparison is not an isomorphism, the
same conclusion for the cover's absolute approximation
does not follow.

For the Fermat tower B.5 applies to the virtually
abelian envelope system, and pro-$\ell$ completion
applies exactly to its $\ell$-power levels. Abelian,
metabelian, bounded nilpotent, and absolute
prounipotent completion have the failures proved
in B.7--B.9. In particular a Selmer point on the
base cannot simply be treated as a section of the
absolute completion of every Fermat cover.
\end{proof}

\begin{statement}{F.12}{Open}{The unrestricted Fermat section sets}
The entire unrestricted section sets for the exact
Fermat approximation systems, and their integral
metabelian quotients, have not been evaluated.
The abelian groups in F.4 and the rational
mixed-Tate Selmer towers in F.2 are computed,
but neither is the requested full set of
splittings of (371)--(373).
\end{statement}
Here is the precise residual problem and the attack made.
After fixing an arithmetic action and a reference section,
one must compute
\[
 H^1_{\rm cts}(G_k,P_{\rm met}),\qquad
 H^1_{\rm cts}(G_k,P^{(\ell)}),\qquad
 H^1_{\rm cts}(G_k,\mathcal V_{\rm ar}),
\]
with the field, packet, and finite-cover descent conditions
of F.11; $\mathcal V_{\rm ar}$ denotes the chosen arithmetic
realization of the virtually abelian envelope system.
The bare pro-discrete group $V(F)$ by itself does not
carry a specified action of all of $G_k$: cyclotomic
scalars already require profinite or adelic coefficient
groups. Its arithmetic realization must therefore be
part of the data, as in B.1, before this expression is
defined.

For the metabelian quotient, the first extension has
kernel the completed Alexander module, acted on
nontrivially through abelianization. The calculation of
that module in B.23 is explicit, but its twisted global
$H^1$ and obstruction $H^2$ for every abelian cocycle
are not evaluated here. Replacing it by successive
rational Tate layers computes F.2 only with its
finite-ramification and crystalline conditions; it
discards integral torsion and full finite monodromy.
For the pro-$\ell$ tower, F.10 gives every finite-stage
obstruction but does not show that all compatible
finite-stage conjugacy classes have compatible
representatives; an inverse-limit argument must also
control their stabilizers. The attempted passage from
the finite-cover comparison to all sections stops at
these explicit steps. F.4 also shows why separation
of cusps cannot be inferred from the abelian calculation.
No unsolved cohomology group in this paragraph is
presented as a computed set.

\begin{statement}{F.13}{Proved}{What the normalizer version does and does not supply}
Replacing the arithmetic image by a Teichm\"uller
normalizer does not by itself define a Selmer problem.
Suppose an action on $U$ extends from an algebraic
arithmetic image $A$ to a proalgebraic normalizer $N$.
Then restriction defines
\[
 H^1_{\rm alg}(N,U)\longrightarrow H^1_{\rm alg}(A,U).
\]
For a faithfully flat image map $G\to A$ through which
the action factors, inflation is injective and its
image consists precisely of cocycles trivial on
$\ker(G\to A)$. Neither operation supplies
unramifiedness or crystallinity at a place.
\end{statement}
\begin{proof}
Restriction is substitution along $A\to N$.
For inflation, two inflated cocycles that differ by
a $U$-coboundary already differ by the same coboundary
on $A$, by faithful flat descent. Conversely a
cocycle trivial on the kernel is constant on its
cosets; descent gives the cocycle on $A$. This is
the nonabelian inflation argument with every map
specified. A local condition refers to a decomposition
group, inertia, or the $p$-adic comparison functors.
An abstract normalizer and its action contain none
of these choices. Consequently the normalizer
calculation in B gives symmetry and compatibility
constraints and possible transports such as F.5;
it cannot replace $G_k$ in the localization diagram
of F.3 without further arithmetic data.
\end{proof}

\begin{statement}{F.14}{Proved}{Lemmas that can be passed to the point-section problem}
The following conclusions apply to a proposed
point-section comparison only with their hypotheses.
A point or rational tangential point gives a
compatible section at every characteristic quotient.
It restricts to a chosen finite cover exactly under
F.11. If a unipotent quotient with the actual
arithmetic action satisfies the local hypotheses
of F.3 or F.8 and a strict global--local dimension
inequality, its integral-point locus is finite.
These facts do not prove that every section is
geometric, or that a finite-level Kummer map is
injective on all point and cusp packets.
\end{statement}
\begin{proof}
Functoriality of the étale fundamental group gives
the first assertion and compatibility of its
quotients. The cover assertion is F.11.
The finiteness assertion is exactly the localization
and analytic argument of F.3, with the refined
elliptic local conditions when F.8 is used. F.4
is an explicit counterexample to finite-level
cuspidal injectivity. The reverse implication
from arbitrary sections to points is not part
of any of these arguments. Thus these are
proved lemmas with the specified action,
base point, local conditions, and dimension
hypotheses; they are not an assumption of
the section conjecture in LM6 Part E.
\end{proof}

